\paragraph{Overview.}
Three findings emerge. (i) A game-state-only model is the best-calibrated wicket
model; adding batter/bowler identity \emph{worsens} held-out log-loss and does not
lower runs RMSE, even though SHAP attributes most explanatory mass to identity.
(ii) Pressure history adds significant skill to expected runs but nothing to
wicket probability. (iii) The apparent death-over momentum effect is largely a
Miller--Sanjurjo finite-sequence artifact; no cell survives correction.

\subsection{Calibrated per-ball benchmark}
\label{sec:results-main}
Table~\ref{tab:wicket} and Table~\ref{tab:runs} report the model ladder, mean over
five seeds. For the wicket head, the state-only model M\_S achieves the best
log-loss ($0.27381$), improving only marginally on the base rate ($0.27484$);
\emph{every} model that adds identity is worse, with the full-identity M\_SBW at
$0.27495$ --- above the base rate. Death-over wickets are, on this evidence,
close to the limit of what the observed state can predict. For the runs head, state captures most of the gain over the
base rate (RMSE $1.728 \to 1.688$); adding identity \emph{raises} RMSE slightly
while \emph{lowering} MAE, indicating that identity helps the typical ball but
adds variance on the extremes. The full model M\_SBW+P+H attains the lowest RMSE
($1.686$) and MAE ($1.233$); the history group, not identity, is what moves RMSE
back below state-only. Seed stability is high (M\_S wicket log-loss std
$7.9\times10^{-5}$; M\_S runs RMSE std $1.4\times10^{-4}$), whereas identity models
show larger seed variance (M\_SBW log-loss std $1.4\times10^{-3}$), consistent
with identity fitting noise. Figure~\ref{fig:ladder} visualizes the ladder.

\begin{table}[t]
\centering
\small
\begin{tabular}{lcccc}
\toprule
Model & log-loss $\downarrow$ & Brier $\downarrow$ & ECE $\downarrow$ & PR-AUC $\uparrow$ \\
\midrule
M0 (base rate)   & 0.27484 & 0.07225 & 0.0000 & 0.0784 \\
\textbf{M\_S (state-only)} & \textbf{0.27381} & 0.07211 & 0.0029 & 0.0904 \\
M\_SB (+batter)  & 0.27413 & 0.07214 & 0.0029 & 0.0880 \\
M\_SW (+bowler)  & 0.27462 & 0.07212 & 0.0019 & 0.0881 \\
M\_SBW (+both)   & 0.27495 & 0.07220 & 0.0026 & 0.0847 \\
M\_SBW+P         & 0.27423 & 0.07216 & 0.0021 & 0.0858 \\
M\_SBW+P+H       & 0.27408 & 0.07214 & 0.0019 & 0.0869 \\
\bottomrule
\end{tabular}
\caption{Wicket head $P(\text{wicket})$ on the held-out 2026 season, mean over
five seeds. State-only (M\_S) is the best-calibrated model; adding identity
worsens log-loss, and no model meaningfully beats the base rate --- death-over
wickets are barely improved upon by any model considered here.}
\label{tab:wicket}
\end{table}

\begin{table}[t]
\centering
\small
\begin{tabular}{lccc}
\toprule
Model & RMSE $\downarrow$ & MAE $\downarrow$ & CRPS $\downarrow$ \\
\midrule
M0 (base rate)   & 1.72835 & 1.29117 & --- \\
M\_S (state-only) & 1.68751 & 1.25512 & --- \\
M\_SB (+batter)  & 1.69154 & 1.24305 & --- \\
M\_SW (+bowler)  & 1.69029 & 1.24746 & --- \\
M\_SBW (+both)   & 1.69208 & 1.24063 & --- \\
M\_SBW+P         & 1.69173 & 1.23991 & --- \\
\textbf{M\_SBW+P+H} & \textbf{1.68639} & \textbf{1.23287} & --- \\
M\_mult (multinomial) & 1.69063 & 1.23843 & 0.81604 \\
\bottomrule
\end{tabular}
\caption{Runs head $E[\text{runs}]$ on the held-out 2026 season, mean over five
seeds. State captures most of the gain over the base rate; identity raises RMSE
but lowers MAE. The full model (state $+$ identity $+$ pressure index $+$ history)
attains the lowest RMSE and MAE, driven by the history group.}
\label{tab:runs}
\end{table}

\begin{figure}[t]
\centering
\includegraphics[width=0.92\linewidth]{figures/fig3_model_ladder.pdf}
\caption{The nested model ladder for both targets. Adding batter/bowler identity
fails to improve held-out wicket log-loss (left) and does not lower runs RMSE
(right); the gains over the base rate come from game state, and only the pressure
history group pushes runs RMSE below state-only.}
\label{fig:ladder}
\end{figure}

\subsection{H1: state versus identity, and the SHAP--skill dissociation}
\label{sec:results-h1}
TreeSHAP attributes the majority of explanatory mass to batter identity ---
$50.6\%$ for runs and $63.7\%$ for wickets --- with state at $31.1\%$/$26.0\%$ and
bowler at $18.3\%$/$10.3\%$ (Figure~\ref{fig:shap}). Yet the held-out incremental
skill of identity over state is zero or negative
(Table~\ref{tab:h1incr}, Figure~\ref{fig:forest}). For the wicket head, adding
batter identity changes log-loss by $-0.00033$ with a confidence interval that
includes zero; for the runs head, adding identity is significantly
\emph{negative} ($\text{B}\mid\text{S}$ RMSE $-0.0053$, 95\% CI $[-0.0089,
-0.0019]$) --- identity hurts out-of-sample RMSE. A leakage audit
(Table~\ref{tab:leakage}) regressing each identity encoding on the state group
finds only modest shared variance (maximum $R^2 = 0.17$ for $\texttt{bat\_deathSR}$),
indicating the encodings are not simply restating state. This rules out the
crudest form of contamination but does not by itself establish the mechanism; the
random-player placebo in Section~\ref{sec:results-robust} provides a more direct
control and indicates that the dissociation reflects in-sample overfitting
of non-stationary player form --- empirical-Bayes historical rates explain training
variance, and hence carry SHAP mass, but do not transfer to the next season's
deliveries.

The predicted asymmetry holds. Identity's share of total state-plus-identity skill
is $-0.137$ for runs and $-0.493$ for wickets, an asymmetry (runs minus wicket)
of $+0.356$: identity is relatively less harmful for runs than for wickets, the
predicted direction. State explains roughly $40\times$ more skill for runs (RMSE
gain $0.0409$ over the base rate) than for wickets (log-loss gain $0.0010$).

\begin{figure}[t]
\centering
\includegraphics[width=0.92\linewidth]{figures/fig1_grouped_shap.pdf}
\caption{Grouped SHAP mass (\% of total $|\text{SHAP}|$) for the full identity
model M\_SBW. Batter identity dominates attribution for both targets ($50.6\%$ for
runs, $63.7\%$ for wickets), despite adding no positive held-out skill
(cf.\ Table~\ref{tab:h1incr}).}
\label{fig:shap}
\end{figure}

\begin{table}[t]
\centering
\small
\begin{tabular}{llrc}
\toprule
Target & Added & Incr.\ skill & 95\% CI \\
\midrule
wicket (log-loss) & B\,$\mid$\,S  & $-0.00033$ & $[-0.00097, +0.00036]$ \\
wicket (log-loss) & W\,$\mid$\,S  & $+0.00024$ & $[-0.00035, +0.00082]$ \\
wicket (log-loss) & BW\,$\mid$\,S & $-0.00033$ & $[-0.00128, +0.00062]$ \\
runs (RMSE)       & B\,$\mid$\,S  & $-0.00531$ & $[-0.00889, -0.00192]$ \\
runs (RMSE)       & W\,$\mid$\,S  & $-0.00271$ & $[-0.00496, -0.00056]$ \\
runs (RMSE)       & BW\,$\mid$\,S & $-0.00492$ & $[-0.00840, -0.00153]$ \\
\bottomrule
\end{tabular}
\caption{Held-out incremental skill of identity over state (match-clustered
bootstrap, 1000 replicates). Positive denotes improvement. For wickets the
intervals include zero; for runs adding identity is significantly negative.}
\label{tab:h1incr}
\end{table}

\begin{figure}[t]
\centering
\includegraphics[width=0.92\linewidth]{figures/fig2_h1_incremental_forest.pdf}
\caption{Forest plot of identity's held-out incremental skill over state, with
match-clustered bootstrap 95\% intervals. No identity addition yields a
significant positive improvement; for the runs head the effect is significantly
negative.}
\label{fig:forest}
\end{figure}

\begin{table}[t]
\centering
\small
\begin{tabular}{lc}
\toprule
Identity feature & $R^2$ vs.\ state \\
\midrule
bat\_deathSR              & 0.169 \\
bat\_boundary\_pct        & 0.151 \\
bat\_dot\_pct             & 0.147 \\
bat\_n\_prior\_balls      & 0.111 \\
bowl\_boundary\_conc\_pct & 0.083 \\
bowl\_n\_prior\_balls     & 0.077 \\
bowl\_deathER             & 0.073 \\
bowl\_dot\_pct            & 0.039 \\
bat\_dismissal\_rate      & 0.020 \\
bowl\_wicket\_rate        & 0.009 \\
\bottomrule
\end{tabular}
\caption{Leakage audit: $R^2$ of each identity encoding regressed on the state
group. Shared variance is modest (maximum $0.17$), so the identity features are
not simply restating state; the random-player placebo
(Section~\ref{sec:results-robust}) provides a more direct control for the
SHAP--skill dissociation.}
\label{tab:leakage}
\end{table}

\subsection{Robustness: a placebo and a six-season rolling window}
\label{sec:results-robust}
The H1 dissociation invites two objections: that the SHAP mass is an artifact of
feeding the model high-variance numeric features rather than real player signal,
and that ``identity hurts'' is specific to the single 2026 test season. We address
both directly.

\paragraph{Random-player placebo.}
We refit the full M\_SBW model with the player-to-encoding mapping randomly
shuffled, so that each player receives another player's historical encoding: the
marginal feature distribution is preserved but player-specific signal is destroyed.
If the SHAP mass merely reflected high-variance features, the placebo would inherit
it; if real identity carried generalizable signal, the placebo would lose skill.
Neither holds (Table~\ref{tab:placebo}, Figure~\ref{fig:placebo}). Batter SHAP mass
collapses from $50.6\%/63.7\%$ (runs/wicket) to $7.5\%/6.9\%$ under the placebo,
with the freed mass flowing to state --- so the attribution is \emph{faithful}: the
model genuinely keys on real player identity. Yet real identity's out-of-sample
skill is actually \emph{worse} than the placebo's: real identity significantly hurts
expected-runs RMSE ($-0.0049$, 95\% CI $[-0.0085, -0.0013]$) while random encodings
are harmless ($+0.0002$, CI spanning zero). The dissociation is therefore genuine
overfitting of non-stationary form, not a leakage or cardinality artifact.

\paragraph{Rolling-window evaluation.}
We repeat the held-out-season protocol for every feasible test season 2021--2026
(train $\le T{-}2$, validate on $T{-}1$, test on $T$). The pattern is stable
(Table~\ref{tab:rolling}, Figure~\ref{fig:rolling}): in \emph{all six} seasons a
state-only model beats state-plus-identity on expected-runs RMSE, and the identity
increment is significantly negative (interval excludes zero) in every season,
ranging from $-0.0049$ to $-0.0131$. Identity never improves wicket log-loss either
(negative point estimate in all six seasons). The ``game state is not improved upon
by identity'' finding is thus a consistent multi-season property, not an artifact of
the 2026 test split.

\begin{table}[t]
\centering
\small
\begin{tabular}{lccc}
\toprule
Encoding & Batter SHAP (runs/wkt) & Runs incr.\ skill (BW$\mid$S) & Wicket incr.\ (BW$\mid$S) \\
\midrule
Real identity   & $50.6\%$ / $63.7\%$ & $-0.0049$ $[-0.0085, -0.0013]$ & $-0.00033$ $[-0.0013, +0.0006]$ \\
Shuffled placebo & $7.5\%$ / $6.9\%$  & $+0.0002$ $[-0.0006, +0.0011]$ & $+0.00008$ $[-0.0004, +0.0005]$ \\
\bottomrule
\end{tabular}
\caption{Random-player placebo. Real identity dominates SHAP attribution yet hurts
held-out runs RMSE; shuffling the player map collapses the attribution and leaves
skill at zero. The faithful-attribution-but-no-skill pattern confirms overfitting
of non-stationary player form.}
\label{tab:placebo}
\end{table}

\begin{table}[t]
\centering
\small
\begin{tabular}{lrccc}
\toprule
Test season & $n$ & Runs RMSE: state $\to$ state$+$id & Identity incr.\ (runs) & Sig.? \\
\midrule
2021 & 14{,}437 & $1.710 \to 1.717$ & $-0.0074$ $[-0.0126, -0.0021]$ & yes \\
2022 & 25{,}545 & $1.701 \to 1.712$ & $-0.0112$ $[-0.0152, -0.0071]$ & yes \\
2023 & 20{,}373 & $1.679 \to 1.692$ & $-0.0131$ $[-0.0183, -0.0082]$ & yes \\
2024 & 28{,}969 & $1.684 \to 1.689$ & $-0.0049$ $[-0.0080, -0.0012]$ & yes \\
2025 & 28{,}202 & $1.700 \to 1.708$ & $-0.0079$ $[-0.0118, -0.0043]$ & yes \\
2026 & 22{,}883 & $1.687 \to 1.692$ & $-0.0049$ $[-0.0085, -0.0013]$ & yes \\
\bottomrule
\end{tabular}
\caption{Rolling-window evaluation. In every held-out season, adding identity
significantly raises expected-runs RMSE (incremental skill negative, interval
excluding zero); state-only is never improved upon. The identity result is not
specific to 2026.}
\label{tab:rolling}
\end{table}

\begin{figure}[t]
\centering
\includegraphics[width=0.95\linewidth]{figures/fig6_identity_placebo.pdf}
\caption{Random-player placebo. Left: batter-group SHAP mass collapses from
$\sim$$50$--$64\%$ to $\sim$$7$--$8\%$ when player identities are shuffled, so the
attribution genuinely tracks real players. Right: real identity significantly hurts
expected-runs RMSE while the placebo is harmless --- faithful attribution, no
generalizable skill.}
\label{fig:placebo}
\end{figure}

\begin{figure}[t]
\centering
\includegraphics[width=0.95\linewidth]{figures/fig7_rolling_window.pdf}
\caption{Six-season rolling-window evaluation. Left: identity's incremental skill on
expected runs is negative with intervals excluding zero in all six held-out seasons.
Right: state-only RMSE is below state-plus-identity in every season. The identity
finding is stable, not a 2026 artifact.}
\label{fig:rolling}
\end{figure}

\subsection{H2: a Markov violation for runs, not wickets}
\label{sec:results-h2}
Adding the pressure-history group on top of the full state-plus-identity-plus-PI
model (M\_SBW+P+H $-$ M\_SBW+P) yields significant incremental skill for expected
runs (RMSE $\Delta = +0.0065$, 95\% CI $[0.0048, 0.0083]$, excluding zero), and
the $\texttt{dot\_run\_len}$ SHAP sign is \emph{negative} --- longer preceding dot
runs lower the next ball's expected runs, a genuine path-dependence
(Table~\ref{tab:h2}, Figure~\ref{fig:h2}). For wicket probability the history
group adds nothing significant (both log-loss and Brier intervals include zero).
The Markov property is therefore violated for scoring intensity but holds for
dismissal risk: the folk claim that ``pressure builds until a wicket falls'' is
not supported once state is controlled; the residual pressure effect is scoring
suppression, not added wicket risk.

\begin{table}[t]
\centering
\small
\begin{tabular}{llrcc}
\toprule
Target & Metric & $\Delta$ & 95\% CI & sig. \\
\midrule
wicket & log-loss & $+0.000165$ & $[-0.000132, +0.000440]$ & no \\
wicket & Brier    & $+0.000021$ & $[-0.000023, +0.000065]$ & no \\
runs   & RMSE     & $\mathbf{+0.006500}$ & $\mathbf{[+0.00481, +0.00826]}$ & \textbf{yes} \\
\bottomrule
\end{tabular}
\caption{H2 incremental skill of pressure history beyond state $+$ identity $+$
pressure index (match-clustered bootstrap). History helps expected runs
significantly (with negative dot-run sign) but adds nothing to wicket probability.}
\label{tab:h2}
\end{table}

\begin{figure}[t]
\centering
\includegraphics[width=0.92\linewidth]{figures/fig4_h2_incremental.pdf}
\caption{Incremental skill of the pressure-history group beyond a full
state-plus-identity-plus-pressure-index model. The effect is significant for
expected runs (interval excludes zero) and null for wicket probability.}
\label{fig:h2}
\end{figure}

\subsection{H3: momentum is a Miller--Sanjurjo artifact}
\label{sec:results-h3}
On the preregistered primary test ($k=1$, scoring shot, hot hand,
state-stratified null), the naive estimator reports a strong negative effect,
$D_{\text{naive}} = -0.114$ --- an apparent ``cold hand'' in which a batter who
just scored is \emph{less} likely to score next. But the Miller--Sanjurjo
finite-sequence bias is $\mathbb{E}_0 = -0.121$, nearly equal and opposite. The
bias-corrected estimate is $D_{\text{corr}} = +0.007$, and the gap is
$D_{\text{corr}} - D_{\text{naive}} = +0.121$ --- the correction moves the answer
by about twelve percentage points, almost entirely accounting for the naive
result. Against the state-stratified permutation null the corrected estimate is
not significant ($p = 0.402$, Benjamini--Hochberg-adjusted $0.9995$).
Figure~\ref{fig:h3} contrasts the naive and corrected estimates across the sweep.

Across the full sweep ($k \in \{1,2,3\}$, two success definitions, two
directions; $1{,}760$ sequences for the primary cell), the naive estimate is
always sizeable ($|D_{\text{naive}}|$ up to $0.31$ at $k=3$) while the corrected
estimate stays small ($|D_{\text{corr}}| \le 0.07$), and the bias grows
monotonically with $k$ --- the signature Miller--Sanjurjo pattern. After
Benjamini--Hochberg correction, \emph{none} of the twelve tests is significant at
$0.05$ (smallest adjusted $p = 0.072$, for boundary/cold/$k=1$, whose raw
state-stratified $p = 0.006$). The naive uncorrected $p$-values flag several
``significant'' cells (e.g.\ $k=3$ raw $p \approx 0.0005$) that are bias
artifacts --- the false-positive pattern the correction is designed to catch. We
therefore find no statistically significant per-ball momentum once the bias is
removed and game state is controlled. We report the $+0.121$ gap as the primary
quantity: an estimate of the apparent momentum introduced by the naive estimator.

\begin{table}[t]
\centering
\small
\begin{tabular}{lrrrrr}
\toprule
Test ($k$, def, dir) & $D_{\text{naive}}$ & bias & $D_{\text{corr}}$ & $p_{\text{strat}}$ & BH \\
\midrule
1, score, hot$^\star$ & $-0.114$ & $-0.121$ & $+0.007$ & 0.402 & 0.9995 \\
2, score, hot & $-0.156$ & $-0.225$ & $+0.069$ & 0.096 & 0.579 \\
3, score, hot & $-0.222$ & $-0.290$ & $+0.069$ & 0.737 & 0.9995 \\
1, score, cold & $+0.090$ & $+0.088$ & $+0.002$ & 0.696 & 0.9995 \\
1, bndry, hot & $-0.131$ & $-0.127$ & $-0.004$ & 0.698 & 0.9995 \\
1, bndry, cold & $+0.038$ & $+0.031$ & $+0.007$ & 0.006 & 0.072 \\
\bottomrule
\end{tabular}
\caption{Selected H3 results on batter death-over stays (full sweep in the
supplementary CSV). $^\star$ denotes the preregistered primary test. The naive
estimate is large in every cell; after removing the Miller--Sanjurjo bias the
corrected estimate is near zero, and no cell survives Benjamini--Hochberg
correction.}
\label{tab:h3}
\end{table}

\begin{figure}[t]
\centering
\includegraphics[width=0.92\linewidth]{figures/fig5_h3_momentum.pdf}
\caption{Naive versus Miller--Sanjurjo-corrected momentum estimates by streak
length. The naive estimate (apparent cold hand) grows with $k$, tracked almost
exactly by the finite-sequence selection bias; the corrected estimate stays near
zero. The $+0.121$ gap at the primary test ($k=1$, scoring shot) is the size of
the momentum illusion.}
\label{fig:h3}
\end{figure}
